\pagestyle{recap}
\stepcounter{section}
\begin{landscape}
\phantomsection\addcontentsline{toc}{section}{Récapitulatif des indexations}
\hypertarget{synthesesujets}{}

\begin{scriptsize}

\section*{Récapitulatif des indexations}

\DTLloaddb{sujetsdb}{sujetsdata.csv}

\rowcolors{1}{white}{gray!25} %

\begin{longtable}{|ccp{1.5cm}c|p{2cm}p{2.7cm}|cp{3.25cm}c|cp{3.25cm}c|}
\hline
\textbf{No} &  \textbf{An.} & \textbf{Lieu} & \textbf{J} & \textbf{Auteur} & \textbf{Titre} & \textbf{Int.} & \textbf{Question d'interpr.} & \textbf{Thème int.} & \textbf{Essai} & \textbf{Question d'essai} & \textbf{Thème ess.}\\\hline
\endfirsthead

\hline
\textbf{No} &  \textbf{An.} & \textbf{Lieu} & \textbf{J} & \textbf{Auteur} & \textbf{Titre} & \textbf{Int.} & \textbf{Question d'interpr.} & \textbf{Thème int.} & \textbf{Essai} & \textbf{Question d'essai} & \textbf{Thème ess.}\\
\hline
\endhead

\hline
\endfoot

%\textbf{No} &  \textbf{An.} & \textbf{Lieu} & \textbf{J} & \textbf{Auteur} & \textbf{Titre} & \textbf{Interprétation} & \textbf{Thème int.} & \textbf{Essai} & \textbf{Thème ess.}\\
\hline
\endlastfoot

\DTLforeach{sujetsdb}{%
  % Mapper les colonnes CSV aux variables LaTeX
  \dbnum=numero,\dblieu=lieu,\dbannee=annee,\dbjour=jour,%
  \dbauteur=nomauteur,\dbtitre=titre,%
  \dbtqun=typequn,\dbtqdeux=typeqdeux,\dbaxun=axeun,\dbaxdeux=axedeux,%
  \dblien=liencorrige,\dbqun=questun,\dbqdeux=questdeux%
}{%
\DTLiffirstrow{}
\\\textbf{\hyperlink{\dblien}{\dbnum{}}} & \dbannee & \dblieu & \dbjour & {\textsc{\dbauteur}} & \textit{\dbtitre}%
	& \StrLeft{\dbtqun{}}{4}{}. & \dbqun{} & \dbaxun{} & \StrLeft{\dbtqdeux{}}{4}{}. & \dbqdeux{} & \dbaxdeux{}
%	& {\detokenize\expandafter{\dblien}}
}

\end{longtable}
\end{scriptsize}
\end{landscape}